%!TEX root = ../thesis.tex

\section{背景}

近年，ヒューマノイドロボットは労働人口減少への対応として，人間に代わる作業の自動化や災害現場などの危険な環境下における，人間の代替手段として期待されている．
人間との類似性，人間に適した環境での移動能力，人間向けに設計された道具の使用能力などから，移動ロボットの中でも重要な位置を占めています．
ヒューマノイドロボットの歩行制御に関する研究は，線形倒立振子モデル（梶田さん論文）をはじめ多数の手法が提案されており，近年は大きくモデルベース（最適制御）と学習ベースのアプローチが存在している．
ヒューマノイドロボットの深層強化学習ベースの歩行制御においての報酬関数の設計が動作（性能？歩行？）に与える影響やトレードオフがあまり議論されていない．
アルゴリズムやフレームワークの提案が多いよねといいたいと
パラメータ調整めんどくさいからchatGPTでやってもらっているけど，人間がどの調整が良いかわかっていない
経験則的にパラメータを決めている


\begin{figure}[hbtp]
  \centering
 \includegraphics[keepaspectratio, scale=0.8]
      {images/RaspberryPiMouse.png}
 \caption{Example}
 \label{Fig:Example}
\end{figure}

\section{関連研究}
Revisiting Reward Design and Evaluation for Robust Humanoid Standing and Walking
Gait-Conditioned Reinforcement Learning with Multi-Phase Curriculum for Humanoid Locomotion
\section{従来研究と課題}
\section{目的}
本研究は，深層強化学習を用いたヒューマノイドロボットの歩行動作学習を対象とする．本研究は，報酬関数を構成する各パラメータを変更しながら学習を行い，得られた各学習モデルを複数の比較指標によって評価・分析する
ことを通じて，報酬関数のパラメータが学習プロセスや獲得される歩行動作に与える影響を明らかにすることである．

\section{論文構成}
本論文の構成を述べる．本論文は，以下の全O章で構成されている．第1章では，本研究の背景や関連技術，目的を述べた．第2章では，本研究に用いた要素技術を述べる．
第3章では，本研究に用いる問題設定について説明し，学習に使用するフレームワークを確認した．
第4章では，本研究の議論にした内容を説明する．
第5章では，実験をする上での予備実験について述べる．
第6章では，実験やその結果に対する考察を述べる．
第7章では，本研究の成果をまとめ，今後の展望について述べる．

\newpage
